\chapter{REVIEW OF LITERATURE AND STUDIES} \label{cap:chapter2}

This chapter reviews the literature that underpins the Multiple Embedding Steganography (MES) framework proposed in this thesis. It is organised in four movements. Section~\ref{sec:ch2-foundations} establishes the conceptual foundations---digital steganography, the RGB colour model, Pixel Value Differencing (PVD), and edge detection as an adaptive embedding criterion. Section~\ref{sec:ch2-studies} surveys the empirical studies that define the current state of the art in spatial-domain, edge-adaptive, and colour-aware embedding. Section~\ref{sec:ch2-security} reviews steganalysis, because the security of a high-capacity scheme cannot be argued from PSNR alone. Section~\ref{sec:ch2-gap} synthesises the review into the specific research gap that MES addresses.

% =========================================================
\section{Related Literatures} \label{sec:ch2-foundations}
% =========================================================

\subsection{Digital Steganography}
Steganography is the practice of concealing a file, message, image, or video within another carrier so that the very existence of the hidden content is concealed. The word combines the Greek \textit{steganos} ("covered") and \textit{graphein} ("writing"). It is therefore distinct in purpose from cryptography: cryptography makes a message unreadable, whereas steganography makes the act of communication itself unremarkable. In the digital image setting, a steganographic scheme is conventionally judged on three competing axes---\textit{capacity} (how much can be hidden), \textit{imperceptibility} (how little the carrier visibly changes), and \textit{security} (how hard the stego-image is to distinguish statistically from a clean one). Recent surveys emphasise that these three axes form a trade-off triangle rather than independent goals: gains on one axis are normally paid for on another, and a contribution is only meaningful when the cost is stated explicitly \citep{kombrink2025survey, janok2026survey}. This thesis positions itself on the capacity axis while committing to explicit imperceptibility thresholds, and treats security as a reported, measured property rather than an assumed one.

\subsection{RGB Color Model}
The RGB colour model is an additive model in which red, green, and blue light are combined to reproduce a broad array of colours, and a standard digital colour image is represented by these three channels. In an 8-bit-depth image each channel holds values from 0 to 255, so each pixel carries 24 bits across three planes. This structure is what makes colour images attractive carriers: a grayscale scheme that embeds $n$ bits per pixel has, in principle, three times that room available in a colour image of the same dimensions.

Two properties of RGB matter for the present work and are easy to overlook. First, the three channels are not statistically independent: natural images exhibit strong inter-channel correlation, and the residual noise of a channel is partly predictable from its neighbours. Second, the Human Visual System (HVS) is not equally sensitive to the three primaries; sensitivity to blue is markedly lower than to green, which is why blue-channel modification is often tolerated at higher amplitudes. The first property is a security liability and is revisited in Section~\ref{sec:ch2-security}; the second is an opportunity that adaptive schemes---including the one proposed here---exploit.

\subsection{Pixel Value Differencing (PVD)}
PVD exploits the differing sensitivity of the human eye to changes in edge areas versus smooth areas. The algorithm partitions the image into non-overlapping blocks of two consecutive pixels and computes the difference between their values. A small difference indicates a smooth region and only a small payload is embedded there, preserving visual quality; a large difference indicates an edge or textured region and a larger payload may be embedded, because the eye tolerates noise in complex textures. The payload per block is therefore data-dependent rather than fixed, which is precisely what gives PVD a better capacity--imperceptibility balance than uniform LSB substitution. The mechanism also fixes the shape of the problem this thesis works on: because the quantisation range table is derived from the cover itself, any scheme built on PVD must guarantee that the receiver can reconstruct the same ranges from the stego-image, and capacity is consequently a property of the cover image, not a constant chosen by the designer.

\subsection{Edge Detection in Steganography}
Edge detectors---Canny, Sobel, Laplacian, and Laplacian-of-Gaussian (LoG) variants---identify regions of sharp intensity change. In steganography these regions are preferred for heavier embedding because the HVS masks distortion in high-frequency content. The practical difficulty is that the edge map must survive embedding: if the detector is run on the cover and the embedding then displaces pixel values across the detection threshold, the receiver recomputing the map on the stego-image obtains a different classification and extraction fails. The literature resolves this in three ways---restricting embedding to bit-planes that do not affect the detector's input, transmitting the map or its seed as side information, or deriving the map from an invariant statistic. Which strategy a scheme adopts is a design commitment with direct consequences for blind extraction, and it is treated as such in Chapter~\ref{cap:chapter3}.

Recent work has moved beyond single-operator edge maps. \citet{ismail2026edgeadaptive} combine Canny and Sobel to obtain a denser and more stable complexity map, and \citet{patwari2025dilatedlog} apply morphological dilation to a LoG map so that pixels adjacent to an edge---not only the edge pixels themselves---become eligible for a heavier payload. The direction of travel is clear: the detector is no longer used as a binary switch but as a graded estimator of local capacity.

% =========================================================
\section{Related Studies} \label{sec:ch2-studies}
% =========================================================

\subsection{Sequential Embedding Techniques}
Many established RGB steganography methods embed sequentially across colour planes: LSB substitution and most PVD variants process the Red channel, then the Green, then the Blue. The approach is straightforward to implement and to verify, but it introduces what may be termed \textit{inter-channel dependency}. Because the per-channel budget is allocated in order, the distortion already spent on earlier channels constrains what later channels may absorb if colour artefacts are to be avoided, and the payload actually realised in the Blue channel is a residual rather than an independently optimised quantity. A second consequence is procedural: a decoding error in an early channel propagates, since the extraction state for later channels depends on the earlier ones having been recovered correctly.

\subsection{Edge-Adaptive and High-Capacity Spatial Schemes}
The most directly comparable recent results come from edge-adaptive spatial-domain schemes, and they establish the capacity bar that any new proposal must clear. \citet{patwari2025dilatedlog} report payloads up to 3.44 bpp at PSNR above 35 dB and SSIM above 0.96 using dilated LoG edge detection with a tunable edge-to-non-edge allocation ratio. \citet{singh2026edgeopt} treat the Canny hysteresis thresholds themselves as optimisation variables and, with variable-bit LSB allocation, report up to 3.3 bpp with PSNR above 34 dB and SSIM above 0.82. \citet{ismail2026edgeadaptive} take the opposite position on the trade-off, accepting a lower payload of approximately 2.97 bpp in exchange for PSNR above 38 dB, SSIM above 0.97, and demonstrated resistance to chi-square and RS steganalysis.

Three observations follow from this cluster of results and they shape the contribution claimed in this thesis. First, payloads in the 3.0--3.5 bpp band are now attainable and reported, so capacity alone is no longer a sufficient contribution; the 3.0 bpp target adopted here is therefore a floor for admissibility, not the claim itself. Second, these schemes operate predominantly on grayscale benchmarks, so the three-fold headroom of a colour carrier is left unexploited and the inter-channel question is never posed. Third, the methods that report the strongest imperceptibility figures are also the ones that report steganalysis results, which suggests that the field now expects security evidence as a matter of course.

\subsection{Colour-Aware and Multi-Channel Schemes}
A smaller body of work engages the colour structure directly. \citet{aljburi2026emd3d} propose a 3D Exploiting Modification Direction scheme in which the R, G, and B values of a pixel are treated as a single coordinate triple, with a 3D Baker chaotic map encrypting the secret before embedding; operating in the decimal domain rather than on a bitstream also yields a real-time throughput advantage. \citet{li2025colorconversion} approach colour from the robustness side, converting between colour representations so that the embedded signal survives the transformations a stego-image typically undergoes in transit.

These two works usefully delimit the space the present research occupies. The 3DEMD scheme is genuinely colour-aware but \textit{couples} the channels: one modification event is defined over the triple, so the three planes cannot be assigned independent, locally derived capacities. The colour-conversion scheme is oriented toward robustness against active distortion, which this thesis explicitly excludes from scope. Neither treats the three planes as three concurrently loaded carriers whose budgets are set independently by their own local edge statistics---which is the architecture proposed here.

\subsection{Learning-Based Image-in-Image Hiding}
For the specific task of hiding an image inside an image, learning-based methods currently report the highest capacities, and the survey by \citet{janok2026survey} traces the architectural progression from CNN and GAN designs through Invertible Neural Networks to Transformer and diffusion-based models. \citet{alrawashdeh2025hed} are notable here because they retain the edge-adaptive principle while replacing the hand-designed detector: Holistically-Nested Edge Detection produces attention maps whose strength modulates the number of bits embedded per pixel, and the resulting stego-images evade XuNet and YeNet at AUC values of 0.47--0.53, statistically indistinguishable from chance.

Two qualifications keep this body of work adjacent to, rather than in competition with, the present thesis. The reported imperceptibility figures---PSNR in the range of 60--61 dB and SSIM of approximately 0.9995---are obtained at roughly 0.1 bpp, two orders of magnitude below the capacity regime targeted here, so the numbers are not comparable without stating the operating point. More fundamentally, learning-based extraction is approximate: the recovered secret is a reconstruction, not a bit-exact copy. For the applications motivating this research---medical imaging, where a single altered pixel may carry diagnostic weight, and secure communication, where the payload may be encrypted and therefore intolerant of any bit error---exact recovery is a requirement rather than a preference. \citet{janok2026survey} further identify benchmarking fragmentation and narrow robustness evaluation as unresolved infrastructure problems in this literature, which is a further reason to prefer a deterministic spatial-domain construction whose behaviour can be reproduced exactly.

\subsection{The Reference Baseline: Liu and Su (2025)}
\citet{liu2025color} propose a colour image steganography method based on the RGB model and edge detection, and their work is the closest antecedent of the present research. They demonstrate that combining edge detection with PVD substantially improves embedding efficiency in colour carriers, and they establish the operational template this thesis inherits: classify local regions by edge strength, and let that classification govern the per-region payload.

Their processing logic, however, remains largely sequential across the colour planes, so the inter-channel dependency described above is carried over rather than removed. This thesis builds on their finding that edge-guided PVD is the right primitive for colour carriers, and diverges on the architectural question of how the three planes are scheduled---proposing a fully parallel arrangement in which each plane is loaded simultaneously and independently, so that the full 24-bit depth of the RGB image is exploited without one plane's budget being contingent on another's.

% =========================================================
\section{Steganalysis and Security Evaluation} \label{sec:ch2-security}
% =========================================================
Steganalysis is the adversarial counterpart of steganography: the attacker's task is to decide whether a given image carries a payload. The field has moved decisively from hand-crafted feature sets with ensemble classifiers to deep convolutional detectors, and \citet{kombrink2025survey} provide a systematic mapping of embedding approaches to the detection strategies that defeat them. \citet{alrusaini2025robustness} evaluate five architectures---EfficientNet, SRNet, ResNet, Xu-Net, and Yedroudj-Net---and find that detector performance is itself fragile under ordinary image transformations, with EfficientNet the most resilient and Xu-Net and Yedroudj-Net degrading sharply under added noise. A practical implication for scheme designers is that a single detector's verdict is weak evidence; a small panel is needed before an imperceptibility claim can be called a security claim.

For this thesis the critical result is \citet{amrutha2023rgbsteganalysis}, who decompose each RGB channel via empirical mode decomposition and combine the resulting features with colour rich model descriptors to detect stego-images across a range of payloads. Their premise is that embedding independently in separate colour planes \textit{disturbs the natural correlation between those planes}, and that this disturbance is detectable even when each plane considered alone appears statistically ordinary. This is a direct and non-obvious challenge to any parallel architecture, the one proposed here included: the very independence that unlocks capacity is the property a colour-aware detector looks for. It does not invalidate the approach, but it does mean that the security of MES must be argued against colour-aware detectors specifically rather than against channel-agnostic ones, and it motivates treating inter-channel correlation preservation as a design constraint rather than an afterthought.

Evaluation infrastructure has improved accordingly. \citet{andone2026stegbench} introduce StegBench, a dual-branch benchmark dataset for multi-class steganalysis spanning JPEG and PNG carriers, which offers a standardised basis for comparison and partially addresses the benchmarking fragmentation that \citet{janok2026survey} identify as a structural weakness of the field.

% =========================================================
\section{Synthesis and Research Gap} \label{sec:ch2-gap}
% =========================================================
Taken together, the reviewed literature supports four conclusions.

\begin{enumerate}
    \item \textbf{Edge-adaptive allocation is settled practice.} Deriving per-region payload from local edge strength is no longer novel in itself; the open questions concern which detector is used and how the allocation is graded \citep{ismail2026edgeadaptive, patwari2025dilatedlog, singh2026edgeopt}.
    \item \textbf{Capacity near 3.0--3.5 bpp is the current bar.} Recent grayscale schemes already reach this band, so a capacity claim must be accompanied by an architectural explanation of \textit{why} the capacity is available and at what imperceptibility cost.
    \item \textbf{Colour carriers remain under-exploited at the architectural level.} Existing colour-aware schemes either couple the channels into a single modification event \citep{aljburi2026emd3d} or target robustness objectives outside the present scope \citep{li2025colorconversion}. The possibility of three concurrently and independently loaded carriers within one image has not been systematically examined.
    \item \textbf{Security evidence is now expected, and for colour schemes it must be colour-aware.} Because inter-channel correlation is itself a detection cue \citep{amrutha2023rgbsteganalysis}, imperceptibility metrics alone are insufficient evidence for a parallel colour architecture.
\end{enumerate}

The gap this thesis addresses follows directly: there is no systematic treatment of the RGB planes as three parallel, independently budgeted carriers with per-channel edge-adaptive PVD, evaluated at high payload against both perceptual metrics and colour-aware detection. The MES framework developed in Chapter~\ref{cap:chapter3} is the response to that gap---targeting an embedding capacity above 3.0 bpp and PSNR above 30 dB, with SSIM reported alongside, while removing the sequential dependency that constrains the Blue-channel payload in existing colour schemes.
